\section{Capítulo~\ref{sec:motorlearning}. O aprendizado motor}

A regra dos mínimos quadrados e a vantagem de um fio de erro separado são teoremas; a construção do circuito com fio de erro, o enfraquecimento por coincidência, o ajuste do traço ao atraso, o pós-efeito e o retorno espontâneo da habilidade foram medidos; que o circuito inteiro funcione como aprendizado supervisionado e construa um modelo interno é a hipótese principal; os números do modelo de dois processos são cálculos pelo modelo do livro.

{\footnotesize
\setlength{\tabcolsep}{3pt}\renewcommand{\arraystretch}{1.15}
\begin{longtable}{>{\raggedright}p{0.5cm}>{\raggedright}p{4.0cm}>{\raggedright}p{5.6cm}>{\raggedright}p{2.3cm}>{\raggedright\arraybackslash}p{1.7cm}}
\toprule
N.º & Afirmação & Experimento e o que foi medido & Fonte & Status \\
\midrule\endfirsthead
\toprule
N.º & Afirmação & Experimento e o que foi medido & Fonte & Status \\
\midrule\endhead
1 & A regra~\eqref{eq:lms} segue, em média, o gradiente do erro quadrático & Resultado da teoria dos filtros adaptativos: uma correção proporcional à entrada e ao erro da saída é uma descida estocástica pelo gradiente do erro quadrático médio. & \cite{WidrowHoff1960} & teoria \\
2 & Com um fio de erro para cada saída, a velocidade não depende do número de saídas; com um único número comum, cai com o número de neurônios ruidosos (proposição~\ref{prop:teachers}) & Curvas de aprendizado de uma rede linear: o aprendizado por perturbações aleatórias dos neurônios é mais lento que o gradiente exato tantas vezes quantos são os neurônios perturbados. & \cite{Werfel2005} & teoria \\
3 & O circuito com fio de erro funciona como aprendizado supervisionado (proposição~\ref{prop:errorwire}) & Teorias deduzidas da estrutura do circuito. Os experimentos das linhas 7--10 concordam com elas; que o circuito inteiro funcione exatamente assim não está provado. & \cite{Marr1969, Albus1971} & hipótese \\
4 & Camada de expansão: cerca de $7\cdot10^{10}$ pequenos neurônios \nt{GLUT}, mais que em todo o resto do cérebro & Contagem de núcleos numa suspensão de tecido dissolvido: no cerebelo humano há $69\cdot10^9$ neurônios dos $86\cdot10^9$ do cérebro inteiro. Estereologia: $101\cdot10^9$ pequenos neurônios no cerebelo de homens idosos. & \cite{Azevedo2009, Andersen1992cereb} & medido \\
5 & Aumentar a dimensão da entrada torna linear a dependência desejada & Teorema sobre o número de dicotomias: uma soma ponderada de $n$ entradas com limiar realiza uma rotulação arbitrária de cerca de $2n$ vetores em posição geral. Que o cerebelo use exatamente isso é parte da hipótese da linha~3. & \cite{Cover1965} & teoria \\
6 & Cerca de $3\cdot10^7$ neurônios \nt{GABA} de saída, cada um com da ordem de $10^5$ entradas & Estereologia do cerebelo humano: $30.5\cdot10^6$ neurônios de saída. Microscopia eletrônica, rato: cerca de $1.75\cdot10^5$ entradas da camada de expansão por neurônio de saída. O neurotransmissor inibitório dos neurônios de saída está na revisão. & \cite{Andersen1992cereb, NapperHarvey1988, Ito2001} & medido \\
7 & Exatamente um fio de erro por neurônio de saída, cerca de uma vez por segundo, cada vez uma descarga poderosa & Revisão de registros: cada neurônio de saída recebe uma entrada \nt{GLUT} trepadeira, que provoca um disparo complexo com frequência da ordem de $1$~Hz. & \cite{Ito2001} & medido \\
8 & A probabilidade da descarga depende da direção do erro de movimento & Macaco, região oculomotora do cerebelo, adaptação de sacadas: a direção do erro visual define a probabilidade do disparo complexo, e a magnitude, o seu momento; a descarga muda os disparos simples na tentativa seguinte e desloca o movimento ao longo do erro preferido da célula. & \cite{Herzfeld2018} & medido \\
9 & A coincidência da descarga de erro com uma entrada ativa enfraquece a entrada ($-\eta e_i x_j$) & Revisão de experimentos: a estimulação conjunta de uma entrada da camada de expansão e do fio de erro provoca um enfraquecimento de longo prazo dessa entrada; a estimulação de uma só entrada, não. & \cite{Ito2001} & medido \\
10 & O comprimento do traço é ajustado ao atraso do erro: com atraso de $100\ms$, enfraquece mais a entrada que esteve ativa $100\ms$ antes & Fatias e camundongos vivos: na região responsável pelo aprendizado oculomotor, o enfraquecimento é sintonizado estreitamente num intervalo de cerca de $120\ms$ entre entrada e erro, o tempo que o sinal de erro leva nessa tarefa; em outra região, células diferentes são sintonizadas em intervalos diferentes. & \cite{Suvrathan2016} & medido \\
11 & O circuito é um filtro adaptativo~\eqref{eq:adaptivefilter} que aprende o modelo interno do corpo & Modelo deduzido da estrutura do circuito. Não há experimento direto em que o filtro aprendido tenha sido medido como função de transferência do corpo. & \cite{Fujita1982} & hipótese \\
12 & A predição subtrai as consequências dos próprios movimentos, por isso não conseguimos fazer cócegas em nós mesmos & Humanos: o próprio toque parece menos cócegas que o mesmo toque alheio; na ressonância, o córtex somatossensorial responde mais fraco ao próprio toque, e o cerebelo fica menos ativo quando o movimento toca a pele. A explicação pela subtração da predição é hipótese. & \cite{Blakemore1998} & hipótese \\
13 & Com uma força externa, o erro diminui e, após a sua retirada, a mão erra para o outro lado & Pessoas movem a manopla de um robô num campo de forças: as trajetórias voltam a ser retas; depois da retirada do campo, o pós-efeito é quase um espelho das primeiras distorções; ele se transfere também para partes não treinadas do espaço. & \cite{ShadmehrMussaIvaldi1994} & medido \\
14 & Aprendem dois processos~\eqref{eq:tworate}; depois da perturbação inversa, a habilidade volta sozinha & Humanos, campo de forças, depois um campo inverso curto e tentativas com o erro travado: a correção volta sozinha à antiga; o modelo de dois processos reproduz isso, o de um só, não. & \cite{Smith2006} & medido \\
15 & Números da fig.~\ref{fig:tworate}: $0.84$ (lento $0.63$) em $200$ movimentos; $6$ inversos; rápido $-0.53$, lento $0.46$; retorno até $0.37$ em $40$ movimentos & Cálculo com \texttt{models/\allowbreak motor\_\allowbreak learning.py}, $A_f=0.92$, $B_f=0.1$, $A_s=0.996$, $B_s=0.02$: $0.840$ ($0.634$ e $0.205$), $6$ movimentos, $-0.531$ e $0.457$, pico de $0.371$ no $40$º movimento. & \texttt{models/\allowbreak motor\_\allowbreak learning.py} & modelo \\
16 & Dois processos são necessários porque o corpo muda em velocidades diferentes & Teoria: a estimativa ótima com perturbações de vários tempos de vida dá vários filtros com constantes de tempo diferentes e explica o retorno espontâneo e o reaprendizado acelerado. & \cite{Kording2007} & hipótese \\
\bottomrule
\end{longtable}}
